\documentclass[11pt]{article}
\usepackage[margin=1in]{geometry}
\usepackage{amsmath,amssymb,amsthm,hyperref}
\newtheorem{lemma}{Lemma}
\newtheorem{theorem}[lemma]{Theorem}
\newtheorem{corollary}[lemma]{Corollary}
\title{A thin-digit terminal family for ordinary gap six}
\author{Internal proof; no priority or external-review claim}\date{17 November 2026}
\begin{document}\setlength{\emergencystretch}{1em}\maketitle
Fix $r\ge1$, $L=7^r$, $N=L+1$, and $c=N+5=L+6$.
Here $N$ is the terminal order.
Put $A=c-zD_z+z$ and
$\Pi_N(X)=\prod_{\ell=0}^{N-1}(X+6\ell)$.
Work over a complete algebraically closed extension of $\mathbb Q_7$,
with valuation ring $\mathcal O$ and residue field $k$.
For $Q\in\mathcal O[X]$, $\deg Q\le4$, define
\[
 G_Q(X)=\binom{X/6+N-1}{N}Q(X),\quad
 d_j(Q)=\sum_{\ell=0}^j(-1)^\ell\binom j\ell G_Q(c-\ell).
\]
Every $d_j$ is integral and $d_j=0$ for $j>N+4$.
Finite formal conjugacy, with no exponential evaluated, gives
\[
 \frac{\Pi_N(A)Q(A)1}{6^N N!}
       =\sum_{j=0}^{N+4}\frac{d_j(Q)}{j!}z^j.             \tag{1}
\]
For $u=z^6=7v$, define
\[
 V_a(Q;v)=\sum_{\ell\ge0}
     \frac{7^\ell a!}{(a+6\ell)!}\,d_{a+6\ell}(Q)v^\ell,
 \qquad0\le a<6.                                       \tag{2}
\]
These six polynomials are integral, since
$v_7(7^\ell a!/(a+6\ell)!)=(s_7(a+6\ell)-a)/6\ge0$.
Their simultaneous vanishing is equivalent to the cyclic zero
of $\Pi_N(A)Q(A)1$ at $u=7v$.

\begin{lemma}[Lucas differences with at most six digits]\label{lem:digit}
For $\bar Q\in k[X]$ put
\[
 f(x)=(x+1)\bar Q(6-x),\qquad
 D_a=\sum_{\ell=0}^a(-1)^\ell\binom a\ell f(\ell).
\]
Let $j\le L+5$ and $s_7(j)=a\le6$.
The only possibly nonzero reduced differences are
\[
 \begin{split}
 \bar d_a&=a(D_a-D_{a-1})\quad(1\le a\le6),\\
 \bar d_{a-1+7^i}&=-D_{a-1}\quad(1\le i\le r,\ 1\le a\le6).
 \end{split}                                                   \tag{3}
\]
All other digit patterns with $s_7(j)\le6$ give zero.
Also $D_6=0$ and $D_5=-[X^4]\bar Q$.
\end{lemma}
\begin{proof}
For an index $\ell$ with $s_7(\ell)\le6$, write $\ell_0$ for
its units digit. The top binomial argument in $G_Q(c-\ell)$ is
\[
 \alpha=1-\ell/6+7^{r+1}/6 .
\]
Since $-1/6=1+7+7^2+\cdots$ in $\mathbb Z_7$, the digits of
$-\ell/6$ through position $r$ are the successive partial sums
of the digits of $\ell$; they have no carry.
Adding1 gives $\alpha_0=1+\ell_0$ and
$\alpha_r=s_7(\ell)$. If $\ell_0=6$, the digit-sum bound forces
$\ell=6$ and the unit digit of $\alpha$ is zero; the formula
below remains zero as well.
Lucas' criterion for $N=7^r+1$ therefore gives
\[
 \overline{G_Q(c-\ell)}
   =(1+\ell_0)s_7(\ell)\bar Q(6-\ell_0)
   =s_7(\ell)f(\ell_0).                                 \tag{4}
\]

In the backward difference defining $d_j$, Lucas' criterion
kills all $\ell$ not digitwise bounded by $j$.
For the others it gives the product of the digit binomials.
The sign $(-1)^\ell$ is $(-1)^{s_7(\ell)}$.
If $j$ has no higher digit, (4) yields
\[
 \sum_{\ell=0}^a(-1)^\ell\binom a\ell\,\ell f(\ell)
                       =a(D_a-D_{a-1}).
\]
If exactly one higher digit is selected and that digit equals1,
its alternating sum kills the $\ell_0$ part of $s_7(\ell)$
and leaves $-D_{a-1}$.
If the total higher digit sum is at least two, its alternating
differences kill the affine function $s_7(\ell)$.
These are all patterns, proving (3), including $j=0$.
Finally $f$ has degree at most five, with leading coefficient
$[X^4]\bar Q$, so $D_6=0$ and
$D_5=(-1)^5 5![X^4]\bar Q=-[X^4]\bar Q$.
\end{proof}

\begin{lemma}[One exact row saturation]\label{lem:matrix}
Put
\[
 E_i=(7^i-1)/6,\qquad
 S(v)=1+\sum_{i=1}^r(-1)^i v^{E_i}.
\]
Every coefficient of $V_0$ is divisible by7.
With columns $Q=1,X,\ldots,X^4$, the reduced rows
$V_1,\ldots,V_5$ form
\[
 M(S)=\begin{pmatrix}
 -1-S&3+S&-S&1+S&4-S\\
 2S&3+S&3&5+5S&3+6S\\
 0&6S&3+6S&5+3S&6+6S\\
 0&0&3S&2+5S&5+6S\\
 0&0&0&S&2+6S
 \end{pmatrix}.                                        \tag{6}
\]
The reduced saturated row $V_0/7$, on arbitrary $Q$, is
\[
                    R(Q)=6Q(6)+vS(v)[X^4]Q.             \tag{7}
\]
\end{lemma}

\begin{proof}
The unit factorial formula is
$7^{-v_7(j!)}j!\equiv(-1)^{v_7(j!)}\prod_i j_i!\pmod7$.
For $1\le a\le5$, only digit sum $a$ contributes to $\bar V_a$.
The low index $j=a$ has weight1. The high index
$j=a-1+7^i$ has normalized unit weight $a(-1)^i v^{E_i}$.
Lemma~\ref{lem:digit} gives exactly
\[
                   \bar V_a=a(D_a-SD_{a-1}).             \tag{5}
\]
Expanding the five basis polynomials $Q$ gives (6).

For $a=0$, every $j>0$ coefficient is7-divisible. Also
\[
 \frac{G_1(c)}7=\frac{1+7L/6}{6(L+1)}
       \prod_{k=1}^{L-1}\left(1-\frac{7L}{6k}\right)
       \equiv1/6=6\pmod7.                               \tag{8}
\]
Every product correction has valuation at least two.
Hence $\overline{d_0(Q)/7}=6\bar Q(6)$.

After division by7, only digit sum six contributes to the
nonconstant part of $V_0/7$.
The low index $j=6$ has reduced difference $6[X^4]\bar Q$
and normalized weight $6v$, giving $v[X^4]\bar Q$.
Each high index $j=5+7^i$ has reduced difference $[X^4]\bar Q$
and weight $(-1)^i v^{E_i+1}$.
Their sum gives $vS(v)[X^4]\bar Q$, proving (7).
\end{proof}


\begin{lemma}[The residual rank locus]\label{lem:rank}
For the six-by-five matrix consisting of (6) and (7), the gcd
of its maximal minors over $\mathbb F_7[v]$ divides $v+1$.
\end{lemma}
\begin{proof}
Here is a direct kernel computation valid for indeterminate $S$.
For $Q$ define
\[
 f(j)=(j+1)Q(6-j),\qquad
 D_a=\sum_{j=0}^a(-1)^j\binom aj f(j),\quad0\le a\le5.
\]
The polynomial $f$ has degree at most five and $f(6)=0$.
The matrix (6), equivalently the operator calculation (3), gives
the finite polynomial identities
\[
 V_a=a(D_a-SD_{a-1}),\quad1\le a\le5,\qquad
                         \sum_{a=0}^5D_a=0.             \tag{9}
\]
The second equality is Newton expansion at $j=6$, using
$\binom6a=(-1)^a$ in $\mathbb F_7$.
The map $Q\mapsto(D_0,\ldots,D_5)$ is an isomorphism onto
this five-dimensional hyperplane.

If $M(S)Q=0$ and $Q\ne0$, then $D_a=S^aD_0$.
Here $D_0\ne0$, so $1+S+\cdots+S^5=0$.
Thus $S=a\in\{2,3,4,5,6\}$ and the kernel is one-dimensional.
Moreover $[X^4]Q=-D_5\ne0$.
Taking the monic kernel representative gives
$D_5=-1$, $D_0=-S$, and hence $6Q(6)=S$.
The additional saturated row (7) is therefore $S(1+v)$.
A rank drop requires $v=-1$.
At $v=0$ one has $S=1$ and $M(S)$ is invertible.

Explicitly,
\[
 \det M(S)=-(1+S+\cdots+S^5).
\]
Its roots in $S$ are simple.
Since $vS'(v)=S(v)-1$, its roots at unit $v$ are simple as well.
The gcd of all maximal minors has no root except possibly $-1$,
and its multiplicity there is at most one because it divides this
determinant. This proves the assertion over $\mathbb F_7[v]$,
not merely at the seven residue-field elements.
\end{proof}

\begin{lemma}[Affine scalar lift]\label{lem:lift}
For fixed $N=7^r+1$, the original characteristic-zero cyclic
matrix has at most one distinct rank-drop point on $u=7v$,
$v\in\mathcal O^\times$.
\end{lemma}

\begin{proof}
The rows (2) are the original cyclic matrix up to nonzero constant
normalizations. Dividing the globally7-divisible row0 by7 preserves
rank at every characteristic-zero point.

Let $\mathcal D_N$ be the determinant of rows $V_1,\ldots,V_5$,
and let $\mathcal E_N$ be the determinant of rows
$V_1,\ldots,V_4,V_0/7$.
Expansion of (6)--(7), with $g(S)=1+S+\cdots+S^5$, gives
\[
 \overline{\mathcal D_N}=-g(S),\qquad
 \overline{\mathcal E_N}=4+3v(S+\cdots+S^5).
\]
Consequently the explicit polynomial
\[
 F_N(v)=2\bigl(\mathcal E_N(v)+3v\mathcal D_N(v)\bigr)
       \in\mathbb Z_7[v],\qquad \bar F_N=v+1              \tag{10}
\]
vanishes at every rank-drop point.
For distinct integral $x,y$, the quotient
$(F_N(x)-F_N(y))/(x-y)$ is a unit: its reduction is1
and all other divided-difference terms are integral multiples of7.
Thus $F_N$ is injective on the entire integral disk, including
ramified extensions, and has at most one zero there.
\end{proof}


\begin{theorem}[A uniform new terminal family]\label{thm:family}
At every terminal order $N=7^r+1$, $r\ge0$, the ordinary factorial
moments of every nonzero complex phase $X(A+BX^6)$ at orders
$N-1,N$ cannot both vanish.
\end{theorem}
\begin{proof}
For $r=0$, $N=2$ is in the already proved Eisenstein family
of the \href{PHYSICAL_PRIME_TOTAL_DEGREE_REDUCTION.tex}{physical-prime reduction},
\texttt{cor:pair}. Now let $r\ge1$ and set
\[
 T(t)=\sum_{k=0}^N\binom Nk(N)_{6k}t^k.
\]
By \texttt{thm:radius} of the same note, any multiple root
satisfies $v_7(u)=1$, $u=-1/t$, in every algebraic embedding.
The \href{gap_three_hurwitz_structure.tex}{all-gap critical-ideal identity},
\texttt{lem:critical-ideal}, identifies such a root with rank
less than five for the cyclic rows $B_N,\ldots,B_{N+4}$.
Since $B_{N+a}=\Pi_N(A)\prod_{\ell=N}^{N+a-1}(A+6\ell)1$,
a unit triangular column change gives precisely the matrix
with columns $Q=1,X,\ldots,X^4$ above.
The imported critical-ideal lemma is proved over $\mathbb C$,
where its exponential gauges are analytic.
For the present fixed integer $N$ and shape $N$, all original
cyclic minors are rational polynomials in $u$.
A multiple root of $T$ is algebraic: embed its number field
into $\mathbb C$, apply that complex critical-ideal implication,
and transport the resulting rational-polynomial minor vanishing
back through any embedding into $\mathbb C_7$.
This establishes the needed rank implication at every algebraic
7-adic root without evaluating any exponential at the boundary.


Lemma~\ref{lem:lift} gives at most one distinct multiple root
of $T$ in $\mathbb C_7$.
If $T$ had a nonconstant rational repeated factor, an irreducible
factor of it would have all its distinct algebraic conjugates
among these same multiple roots. Characteristic-zero separability
would force that irreducible factor to have degree one and hence
a rational real root. But $N$ is even and
\[
 T(t)=\frac1{\Gamma(N)}\int_0^\infty
            x^{N-1}e^{-x}(1+tx^6)^N\,dx>0
 \qquad(t\in\mathbb R).
\]
This is impossible. Thus $T$ is squarefree over $\mathbb Q$
and over $\mathbb C$.

For $AB\ne0$, the normalized consecutive moment polynomials are
$T+(6/N)tT'$ and $T-(1/N)tT'$.
Their coefficient determinant is $-7/N\ne0$ and $T(0)=1$;
a common zero would be a multiple root of$T$.
Nonzero degenerate monomials have no zero positive moment.
\end{proof}

\begin{corollary}[Algebraic Gamma shapes in a terminal-dependent disk]
\label{cor:gamma-disk}
Let $r\ge1$, $N=7^r+1$, and let $\beta>0$ be real algebraic.
Suppose one embedding $\mathbb Q(\beta)\longrightarrow\mathbb C_7$
satisfies $v_7(\beta-1)\ge r+2$.
Then the Gamma-$\beta$ moments of every nonzero complex
$X(A+BX^6)$ at orders $N-1,N$ cannot both vanish.
\end{corollary}
\begin{proof}
The case $\beta=1$ is the ordinary theorem. Otherwise put
$\delta=\beta-1\ne0$ and $\gamma=N+\delta>0$.
The bridge polynomial is $T_{N,\gamma,6}$.
We first establish the same necessary boundary radius in the
possibly ramified coefficient field.
For every integer factor $a$ in $N,\ldots,7N-1$,
$v_7(a)\le r+1$, so $1+\delta/a\in1+7\mathcal O$.
Consequently
\[
 H_\gamma(Y)=7T_{N,\gamma,6}(Y/7)
                  \equiv H_N(Y)\pmod7
\]
coefficientwise, where $H_N=7T_{N,N,6}(Y/7)$.
It is integral, its leading coefficient is a unit,
and its linear coefficient is a unit, because $N=1$ modulo7.
All its roots are integral. At a root of positive valuation
its derivative is a unit, so every multiple root must be a unit.
Thus every hypothetical multiple root of $T_{N,\gamma,6}$
still satisfies $u=7v$, $v\in\mathcal O^\times$.

The actual raising constant is now $c+\delta$.
For $x\in\mathbb Z_7$ and $1\le k\le N$, one has
\[
 v_7\binom{\delta/6}{k}
          =v_7(\delta)-v_7(k)\ge2.
\]
Indeed for $1\le j<k\le N$ one has
$v_7(j)\le r<v_7(\delta/6)$, so
$v_7(\delta/6-j)=v_7(j)$; cancellation in the binomial
product gives the displayed valuation.
Vandermonde's identity therefore gives
$\binom{x+\delta/6}{N}-\binom xN\in49\mathcal O$.
For every integral primitive coefficient polynomial $Q$,
its evaluations at $c+\delta-\ell$ and $c-\ell$ differ by
$\delta\mathcal O\subset49\mathcal O$ as well.
Thus all actual $d_j(Q)$ differ from their ordinary counterparts
by $49\mathcal O$.
The weights in (2) are integral, so every $V_a$ changes by
$49\mathcal O$, and $V_0/7$ changes by $7\mathcal O$.
The same matrices and explicit minor combination (10) now give
$F_{N,\gamma}\in\mathcal O[v]$ with reduction $v+1$.
It is injective on the whole integral disk.

The critical-ideal implication transfers algebraically from
$\mathbb C$ just as in the ordinary proof, now over the real
coefficient field $E=\mathbb Q(\beta)$.
All $E$-conjugates of a repeated irreducible factor, with $\beta$
fixed and at the one chosen7-adic embedding, enter the same
matrix and unit domain. The affine minor permits at most one
distinct root; separability forces a degree-one factor over
$E\subset\mathbb R$. This contradicts strict even-moment
Gamma positivity, since $N$ is even and $\gamma>0$.

Finally, the actual normalized moments are
\[
 T_{N,\gamma,6}+\frac6\gamma\,tT'_{N,\gamma,6},
 \qquad T_{N,\gamma,6}-\frac1N\,tT'_{N,\gamma,6}.
\]
Their determinant is $-(6/\gamma+1/N)\ne0$.
A common root would be a multiple root of $T_{N,\gamma,6}$,
already excluded. Degenerate monomials are immediate.
\end{proof}

\begin{corollary}[The refined degree-seven frontier]
\label{cor:degree-seven-frontier}
Let
\[
 {\cal U}=\{b7^s:1\le b\le6,\ s\ge0\}
                \cup\{7^r+1:r\ge0\},\qquad
 P_j(t)=\sum_{k=0}^j\binom jk(j+1)_{6k}t^k.
\]
Strong Factorial for every ordinary complex univariate binomial
of degree at most seven is equivalent to
\[
 \gcd_{\mathbb Q[t]}(P_{N-1},P_N)=1
              \quad\text{for every }N\ge6,\quad N\notin{\cal U}.
\]
Any hypothetical common root on these remaining tests still
satisfies $v_7(-1/t)=1$ at every algebraic embedding.
These remaining tests, all-index ordinary gap six, and degree-seven
binomial SFC are not proved here.
\end{corollary}
\begin{proof}
Theorem \texttt{thm:gap-six-frontier} of the
\href{ORDINARY_DEGREE_SEVEN_GAP_SIX_FRONTIER.tex}{earlier gap-six frontier}
gives the exact equivalence after omitting the first family in
$\mathcal U$ and supplies the necessary root radius.
Theorem~\ref{thm:family} supplies every terminal in the second
family. Removing these already proved tests preserves sufficiency;
necessity is immediate since the remaining tests are a subset
of the degree-seven assertion.
\end{proof}
The ordinary result covers an unbounded terminal family beyond
$b7^s$, $1\le b\le6$.
The algebraic Gamma corollary uses an order-dependent disk at one
specified place; it is not an all-positive-real-shape assertion.
Neither statement proves all-index ordinary gap six,
degree-seven binomial SFC, or unrestricted FC/SFC.
\end{document}
